\section{Complex analysis with a radius gap and simultaneous derivative control}
\label{sec:complex}

The function-space construction in Section~\ref{sec:function-topology}
requires an explicit common modulus. For analytic families that information
can be produced from a coefficient bound and a radius gap. We give a
direct construction that controls every derivative up to a requested order,
using only rational polynomial calculations and geometric bounds.
This is a quantitative normal-family statement in a \(C^J\) metric, without
asserting the existence of a computable subsequence.

The information supplied with an analytic function is significant.
Representing a function in a continuous-function space and additionally
knowing that it is analytic is not generally equivalent to having an analytic
representation with effective radius and bound data
\cite{PaulySteinberg2018}. We therefore specify this data as part of the
presentation and show exactly what follows from it.

\subsection{A coefficient presentation}

Fix positive rational \(R\), nonnegative rational \(M\), and coefficients
\(a_n\in\C\) with a joint primitive-recursive Gaussian-rational evaluator.
The normalized coefficients
\[
                         b_n=R^n a_n
\]
have Gaussian-rational approximants to every dyadic precision, and the input
certificate asserts
\[
                             |b_n|\leq M\quad(n\in\N).
\tag{A1}
\]
This is a schema in one fixed joint evaluator. It does not enumerate all
analytic functions. The case \(M=0\) yields the zero function and can be
handled separately.

Put \(P_N(z)=\sum_{n=0}^Na_nz^n\).
For \(j\geq0\) write
\((n)_j=n(n-1)\cdots(n-j+1)\), with \((n)_0=1\) and
\((n)_j=0\) when \(n<j\). Then
\[
                 P_N^{(j)}(z)=
                   \sum_{n=j}^N(n)_j a_n z^{n-j}
\]
is a finite identity. The superscript is initially only formal polynomial
differentiation; the theorem below proves agreement with analytic
differentiation.

\begin{lemma}[A finite differentiated geometric bound]
\label{lem:derivative-geometric}
For rational \(0\leq t<1\), all \(j,L\in\N\),
\[
            \sum_{n=j}^L(n)_j t^{n-j}
                      \leq \frac{j!}{(1-t)^{j+1}}.
\tag{A2}
\]
The assertion is a finite rational inequality, with no differentiated
infinite series in its proof.
\end{lemma}
\begin{proof}
If \(L<j\), the left side is zero. Otherwise the following finite expansion
has nonnegative degrees throughout.
There are \(\binom{m+j}{j}\) ordered \((j+1)\)-tuples of nonnegative
integers with sum \(m\). This identity follows by finite induction, or by
placing \(j\) separators among \(m+j\) positions.
Expand
\((1+t+\cdots+t^L)^{j+1}\).
All its coefficients are nonnegative, and its terms of total degree
at most \(L-j\) include every such tuple. Consequently
\[
\sum_{m=0}^{L-j}\binom{m+j}{j}t^m
 \leq(1+t+\cdots+t^L)^{j+1}
 \leq(1-t)^{-j-1}.
\]
Multiply by \(j!\) and use
\((m+j)_j=j!\binom{m+j}{j}\).
If \(L<j\), the left side of (A2) is zero.
The tuple count, expansion, and geometric inequality are all finite
primitive-recursive calculations.
\end{proof}

\subsection{A uniform tail for all requested derivatives}

Fix rational \(0\leq r<R\) and an order \(J\in\N\). Define
\[
s=\frac{r+R}{2},\qquad
\rho=\frac rR,\qquad \theta=\frac rs,\qquad \sigma=\frac sR ,
\]
so \(0\leq\theta<1\), \(0<\sigma<1\). For \(0\leq j\leq J\) put
\[
 K_j=\frac{M j!}
 {s^j(1-\theta)^{j+1}(1-\sigma)},\qquad
 D_j=\frac{j!}{R^j(1-\rho)^{j+1}}.
\tag{A3}
\]
These constants are rational and explicitly computable, including at \(r=0\).

\begin{theorem}[Derivative-tail estimate and holomorphic interpretation]
\label{thm:analytic-tail}
For every \(L\geq N\geq J\), \(0\leq j\leq J\), and \(|z|\leq r\),
\[
          |P_L^{(j)}(z)-P_N^{(j)}(z)|
                          \leq K_j\sigma^{N+1}.
\tag{A4}
\]
The series thus supplies jointly primitive-recursive uniform names of
\(f,f',\ldots,f^{(J)}\) on the smaller disk, and \(f\) is holomorphic on
\(|z|<R\).
\end{theorem}
\begin{proof}
For \(n>N\geq j\), the coefficient bound gives
\[
 |(n)_ja_nz^{n-j}|
 \leq M s^{-j}(n)_j\theta^{n-j}\sigma^n.
\]
By (A2), each nonnegative factor
\((n)_j\theta^{n-j}\) is at most
\(j!/(1-\theta)^{j+1}\).
The finite sum of \(\sigma^n\), \(N<n\leq L\), is at most
\(\sigma^{N+1}/(1-\sigma)\).
Multiplication proves (A4). This also constructs the limit names:
choose \(N\) by the explicit clock below and evaluate the finitely many
coefficients to a precision whose finite weighted error is small enough.
No unbounded search for a convergent tail is involved.

To justify analytic differentiation rather than merely name a formal
derivative, first work with \(j=0,1,2\) and an intermediate disk.
For a finite polynomial, binomial expansion gives
\[
P_N(z+h)-P_N(z)-hP_N'(z)
 =h^2\sum_{n=2}^N a_n
               \sum_{\ell=2}^n\binom n\ell z^{n-\ell}h^{\ell-2}.
\]
If \(|z|\leq r\) and \(|h|\leq s-r\), the absolute value of the right side
is bounded by
\[
\frac{|h|^2}{2}\sum_{n=2}^N n(n-1)|a_n|s^{n-2}
 \leq \frac{M}{R^2(1-s/R)^3}|h|^2 .
\tag{A5}
\]
The first inequality follows termwise from
\(\binom n\ell\leq n(n-1)\binom{n-2}{\ell-2}/2\)
and \(|z|+|h|\leq s\); the second is (A2).
Pass to arbitrarily accurate polynomial approximants in (A5), with the
explicit errors from (A4). The resulting difference-quotient error is at
most a known constant times \(|h|\), proving that the named first
derivative is the complex derivative. Division by \(h\) is used only for
increments apart from zero; the remainder inequality itself holds without
that restriction.
The same argument applied to the formally differentiated polynomials,
using (A2) with the higher orders required, proves every further
derivative. For derivative order \(j\), an explicit remainder constant is
\[
             \frac{M(j+2)!}{2R^{j+2}(1-s/R)^{j+3}}.
\]
For a chosen interior point one chooses a rational smaller disk
and collar before making this argument. Thus holomorphicity and all
differentiation claims follow from finite identities and specified moduli.
\end{proof}

\paragraph{Relation with Cauchy integral presentations.}
One can instead supply a function on a certified neighborhood of
\(\{|z|\leq R\}\), a bound \(M\), a primitive-recursive contour evaluator,
and a proved Cauchy formula. The coefficient formula then gives (A1).
If the bound is available only on smaller circles, include the corresponding
radius loss. The repository's Cauchy and residue packet develops the
required finite contour estimates under explicit differentiability and
geometric hypotheses \cite{repoCauchy}; it is a proposed source of this
adapter, not a silently completed new Lean dependency.
The present theorem itself starts at the coefficient certificate and needs
no contour theorem as an unproved premise.

\subsection{A completely explicit finite \texorpdfstring{\(C^J\)}{C-J} approximation family}

Let
\[
\begin{split}
K&=\max_{j\leq J}K_j,\qquad D=\max_{j\leq J}D_j,\\
U&=\max(1,\lceil K\rceil),\qquad
V=\max(1,\lceil8D\rceil),\qquad L_0=\max(1,\lceil M\rceil),\\
h&=\ceil{1/(1-\sigma)},\\
N&=J+h(k+U+2),\qquad p=k+V,\qquad \Delta=2^{-p}.
\end{split}
\tag{A6}
\]
Use the real rational grid
\[
       G_p=\{j\Delta:j\in\Z,\ -L_0\,2^p\leq j\leq L_0\,2^p\}.
\]
For every \(N+1\) tuple \(q_0,\ldots,q_N\in G_p+iG_p\), form
\[
                         Q_q(z)=\sum_{n=0}^N q_n R^{-n}z^n .
\tag{A7}
\]
The list contains exactly
\[
                      H=(2L_0\,2^p+1)^{\,2(N+1)}
\tag{A8}
\]
polynomials. Listing all of them can be extremely expensive, but is
primitive recursive. A chosen approximant can be computed without
enumerating the entire list.

\begin{theorem}[Simultaneous derivative entropy bound]
\label{thm:analytic-entropy}
For every family member satisfying (A1), finite observations select a
polynomial \(Q_q\) in (A7) for which
\[
 \max_{0\leq j\leq J}\ 
        \sup_{|z|\leq r}|f^{(j)}(z)-Q_q^{(j)}(z)|
                        \leq2^{-k-1}.
\tag{A9}
\]
The selecting algorithm, degree, coefficient precision and cardinality
bound are primitive recursive in the supplied rational data, \(J,k\),
and the fixed coefficient evaluator.
\end{theorem}
\begin{proof}
The proof of Lemma~\ref{lem:clock}, with \(q=\sigma\), gives
\(\sigma^h\leq1/2\). Since \(K\leq U\leq2^U\),
\[
K\sigma^{N+1}\leq U\,2^{-(k+U+2)}\leq2^{-k-2}.
\tag{A10}
\]
This is a common bound for the derivative tails.

Approximate each real and imaginary part of \(b_n\), \(0\leq n\leq N\),
within \(\Delta/4\), round to a nearest grid point, and clip to
\([-L_0,L_0]\). The true component lies in this interval, so clipping cannot
increase its error. Each coordinate error is at most \(3\Delta/4<\Delta\),
and hence \(|b_n-q_n|\leq2\Delta\).
At derivative order \(j\), (A2) gives the uniform polynomial error
\[
\left|\sum_{n=j}^N(n)_j(a_n-q_nR^{-n})z^{n-j}\right|
 \leq 2\Delta R^{-j}\sum_{n=j}^N(n)_j\rho^{n-j}
 \leq2\Delta D.
\]
By \(V\geq8D\) and \(2^V\geq V\), this is at most \(2^{-k-2}\).
Add (A10) to obtain (A9).
All selection decisions are comparisons of rational observations; none
is a decision about equality or sign of an exact complex coefficient.
\end{proof}

The approximating polynomials lie in the ambient \(C^J\) space; they need not
belong to the original analytic family or obey its original coefficient
bound. Selecting representatives inside an arbitrary family would be
an additional existence problem. The theorem gives an explicit total
boundedness result for the analytic family viewed in the smaller-disk
\(C^J\) metric. It does not choose a convergent subsequence from an arbitrary
sequence in that family.

\subsection{An example with exact rational derivative checks}

Take \(R=2\), \(M=1\) and \(a_n=2^{-n}\). Then
\[
                      f(z)=\frac1{1-z/2},\qquad
          f^{(j)}(z)=\frac{j!\,2^{-j}}{(1-z/2)^{j+1}}.
\]
For rational complex \(z\) with \(|z|\leq r<2\), both the finite polynomial
jets and the right-hand expressions are exactly Gaussian rational.
The reference program compares them using squared norms and verifies
(A4), (A10) and (A9)'s error budgets at several radii and derivative orders.
The integer schedules in (A6) are intentionally conservative; their purpose
is an unconditional primitive-recursive bound, not good numerical complexity.

\subsection{The radius gap and the analytic structure cannot be discarded}

On the closed unit disk the polynomials \(z^n/n\) tend uniformly to zero,
while their derivatives \(z^{n-1}\) have uniform norm \(1\).
Thus differentiation is not continuous for the same-disk uniform metric.
The supplied smaller radius changes the conclusion: the powers of
\(\sigma<1\) absorb all polynomial derivative factors. This explains the
appearance of a collar in analytic arguments, rather than treating it as
a technical artifact.

The continuous partitions of unity from the topology section must also be
used with care. They glue approximately compatible continuous data, but
their weights are generally not holomorphic. An approximately glued result
is therefore not automatically holomorphic. If the local holomorphic
functions agree exactly on the relevant overlaps, the glued function equals
them locally and is holomorphic; otherwise a separate analytic correction
argument is needed. This is one reason to retain domain-specific
interpretation lemmas even in a common framework.

\subsection{Object statements and the Lean path}

Encode a Gaussian rational as a pair of rationals, a polynomial as a finite
coefficient list, and a jet as the list of its formal derivatives through
order \(J\). Horner evaluation, finite differentiation, factorials, powers,
grid rounding, and (A6) are fixed primitive-recursive programs.
The core object statements quantify over finite \(L,N,J\), a rational query
\(z\), and rational coefficient observations, with explicit error terms for
those observations. The majorant conclusion is a rational squared-norm
inequality; all denominators \(R,s,1-\rho,1-\theta,1-\sigma\) have proved positive
lower bounds.

The needed derivations are: the finite geometric identity; (A2) by a bounded
coefficient count; a triangle estimate for finite complex sums; the monomial
derivative identity; the split-radius majorant (A4); the dyadic clock;
the rounded-coefficient estimate; and the finite difference identity (A5).
Their induction invariants are bounds for finite partial sums or coefficient
lists, not holomorphicity assertions. After these are obtained, regularity,
extensionality, and analytic interpretation are assembled as in
Theorem~\ref{thm:compilation}. The actual object proof terms for these new
complex-analytic statements remain explicit formalization targets.
